%(BEGIN_QUESTION)
% Copyright 2010, Tony R. Kuphaldt, released under the Creative Commons Attribution License (v 1.0)
% This means you may do almost anything with this work of mine, so long as you give me proper credit

Finn regnearket «{\bf grade template}», og undersøk hvordan du kan beregne din egen karakter i et valgfritt andreårsemne ved å fylle inn poeng for eksamener (både mastery og proporsjonal -- skriv inn et tall større enn 1 for å simulere ny prøve på en mastery-eksamen), quizer, forsentkomming, fravær og andre variabler. Merk forskjellen mellom karakterberegning i dette INST200-emnet og emner med flere studiepoeng (for eksempel INST240, INST250, INST260).

\vskip 10pt

Diskuter betydningen av å signere FERPA-skjema, slik at potensielle arbeidsgivere kan få innsyn i detaljer om dine faglige prestasjoner. Uten signert samtykke kan en instruktør ikke kommentere studenten utover å opplyse om vedkommende er registrert i et emne eller ikke.

\vskip 10pt

Undersøk karakterarket for å identifisere data en arbeidsgiver kan være interessert i ved vurdering av ansettelse. {\it Merk at det registreres mye mer enn bare sluttkarakter eller prosentscore.} Hvilket datapunkt tror du de fleste arbeidsgivere primært fokuserer på?

\vskip 20pt

Finn også kvartalets {\bf calendar}-regneark, og identifiser hvordan kalenderdatoer mappes mot dagene som vises på forsiden av hver seksjonsoppgave.

\vskip 20pt \vbox{\hrule \hbox{\strut \vrule{} {\bf Forslag til sokratisk diskusjon} \vrule} \hrule}

\begin{itemize}
\item{} Hvorfor tror du instruktøren gir deg et regneark for å beregne egen karakter i stedet for å bare fortelle hvilke karakterer du har oppnådd underveis i programmet?
\item{} Hvorfor tror du det finnes et separat kalenderdokument som knytter arbeidsarkseksjoner og dager til kalenderdager? Hvorfor ikke bygge disse datoene direkte inn i PDF-arbeidsarkene (på side 1)?
\end{itemize}

\underbar{file i00118}
%(END_QUESTION)





%(BEGIN_ANSWER)

Hint: både karaktermalen og kalenderregnearket ligger på nettverksdisken {\tt Y:} ved BTC, i Microsoft Excel-format.

%(END_ANSWER)





%(BEGIN_NOTES)


%INDEX% Course organization, grading: spreadsheet-based grade calculator
%INDEX% Course organization, schedule: spreadsheet-based calendar

%(END_NOTES)
